\documentclass[UTF8,a4paper,11pt]{ctexart}
\usepackage[margin=2.35cm]{geometry}
\usepackage{amsmath,amssymb,mathtools,bm}
\usepackage{booktabs,longtable,array}
\usepackage{xcolor}
\usepackage{listings}
\usepackage{enumitem}
\usepackage{hyperref}
\usepackage{fancyhdr}
\usepackage{lastpage}
\usepackage{microtype}

\hypersetup{colorlinks=true,linkcolor=blue!55!black,urlcolor=blue!60!black,citecolor=blue!55!black}
\pagestyle{fancy}
\fancyhf{}
\lhead{第四题 4.1：Poisson 方程}
\rhead{有限体积离散与 OpenFOAM 求解器}
\cfoot{\thepage/\pageref{LastPage}}
\setlength{\headheight}{14pt}
\setlength{\parindent}{2em}
\setlength{\parskip}{0.35em}
\setlength{\emergencystretch}{3em}
\setlist{nosep,leftmargin=2em}

\definecolor{codebg}{RGB}{247,248,250}
\definecolor{codekw}{RGB}{0,74,124}
\definecolor{codestr}{RGB}{145,42,42}
\definecolor{codecom}{RGB}{62,116,62}
\lstset{
  basicstyle=\ttfamily\small,
  backgroundcolor=\color{codebg},
  frame=single,
  rulecolor=\color{black!18},
  breaklines=true,
  columns=fullflexible,
  keepspaces=true,
  keywordstyle=\color{codekw}\bfseries,
  commentstyle=\color{codecom},
  stringstyle=\color{codestr},
  showstringspaces=false,
  numbers=left,
  numberstyle=\tiny\color{black!55},
  numbersep=7pt
}
\lstdefinelanguage{OpenFOAM}{
  morekeywords={FoamFile,version,format,class,object,dimensions,internalField,boundaryField,type,value,default,none,solver,preconditioner,tolerance,relTol,residualControl,nNonOrthogonalCorrectors},
  sensitive=true,
  morecomment=[l]{//},
  morecomment=[s]{/*}{*/},
  morestring=[b]"
}

\newcommand{\vect}[1]{\bm{#1}}
\newcommand{\Sf}{\vect S_f}
\newcommand{\Scf}{\vect S_{cf}}
\newcommand{\df}{\vect d_f}
\newcommand{\nf}{\vect n_f}
\newcommand{\kf}{\vect k_f}
\newcommand{\grad}{\vect\nabla}

\title{第四题第 4.1 问详细解答\\[0.4em]
\large Poisson 方程的有限体积空间离散与 OpenFOAM v14 求解器}
\author{}
\date{}

\begin{document}
\maketitle

\section{题目与必要说明}

控制方程为
\begin{equation}
  \nabla\!\cdot\!\nabla\phi=\omega,
  \qquad\text{即}\qquad
  \nabla^2\phi=\omega .
  \label{eq:pde}
\end{equation}
要求基于有限体积法写出空间离散形式，并采用 OpenFOAM 提供的空间离散格式开发数值求解器。

方程~\eqref{eq:pde} 是稳态 Poisson 方程，本身没有时间导数。因此题目中的“半离散形式”应理解为：连续区域上的偏微分方程只在\textbf{空间上}作有限体积离散，得到单元中心未知量满足的代数方程。严格地说，它不是“保留连续时间变量”的时间半离散方程。

还需指出：Poisson 方程必须配合边界条件才构成闭合问题。第 4.1 问未给出具体边界值，以下先写一般边界通量；随后给出规定值边界
\(\phi=g\)（OpenFOAM 中的 \texttt{fixedValue}）的离散。若全部边界均为 Neumann 条件，则还需满足整体相容条件并固定一个参考值，否则解只确定到一个加法常数。

本文的有限体积推导主要依据 Moukalled、Mangani 与 Darwish 的教材第 8 章，特别是式 (8.4)、(8.63)--(8.72)、(8.75)--(8.77) 以及第 8.10.2 节\cite{Moukalled2016}；OpenFOAM 实现细节按 OpenFOAM Foundation v14 用户指南和相应源代码\cite{OFGuide,OFfvmlap,OFcorrected,OFsurface,OFfixedValue}核对。

\section{统一记号}

全文采用此前两题已经统一的 owner--neighbour 记号，不使用含义不清的相邻单元记号 \(d(f)\)，也不以 \(A_f\) 表示面面积。

\begin{longtable}{>{\centering\arraybackslash}p{0.22\textwidth}p{0.70\textwidth}}
\toprule
记号 & 含义\\
\midrule
\endfirsthead
\toprule
记号 & 含义\\
\midrule
\endhead
\(c,\Omega_c,V_c\) & 当前控制体、控制体区域及其体积\\
\(\mathcal F_c^{\mathrm{int}}\) & 与控制体 \(c\) 相接的内部面集合\\
\(\mathcal F_c^{\mathrm{bnd}}\) & 控制体 \(c\) 上的边界面集合\\
\(O_f,N_f\) & 内部面 \(f\) 的 owner 单元和 neighbour 单元\\
\(\vect x_{O_f},\vect x_{N_f}\) & owner 与 neighbour 的单元中心\\
\(\df=\vect x_{N_f}-\vect x_{O_f}\) & 从 \(O_f\) 指向 \(N_f\) 的中心连线向量\\
\(d_f=\lVert\df\rVert\) & 两单元中心距离\\
\(\nf\) & 从 \(O_f\) 指向 \(N_f\) 的全局单位法向量\\
\(\Sf=|S_f|\nf\) & 按 \(O_f\to N_f\) 定向的全局有向面积向量\\
\(\sigma_{cf}\) & 若 \(c=O_f\) 则为 \(+1\)，若 \(c=N_f\) 则为 \(-1\)\\
\(\Scf=\sigma_{cf}\Sf\) & 相对于单元 \(c\) 的外向面积向量；亦即 \(\Scf=|S_f|\vect n_{cf}\)\\
\(H_f\) & 按全局方向 \(O_f\to N_f\) 定义的内部面扩散通量\\
\(Q_c\) & 尚未除以体积的单元离散残差\\
\(R_c=Q_c/V_c\) & 除以体积后的空间残差\\
\bottomrule
\end{longtable}

这里 \(|S_f|\) 是面的几何面积；二维网格在有限体积意义下可理解为面长度乘挤出厚度。局部和全局面积向量之间始终满足
\begin{equation}
  \Scf=\sigma_{cf}\Sf=|S_f|\vect n_{cf},
  \qquad
  \sigma_{cf}=
  \begin{cases}
  +1,&c=O_f,\\
  -1,&c=N_f.
  \end{cases}
  \label{eq:sigma}
\end{equation}

\section{从积分方程到有限体积空间离散}

\subsection{控制体积分与 Gauss 定理}

在任意控制体 \(\Omega_c\) 上积分方程~\eqref{eq:pde}：
\begin{equation}
  \int_{\Omega_c}\nabla\!\cdot\!(\nabla\phi)\,\mathrm dV
  =\int_{\Omega_c}\omega\,\mathrm dV .
  \label{eq:cvint}
\end{equation}
应用 Gauss 散度定理，得到精确的控制体平衡式
\begin{equation}
  \oint_{\partial\Omega_c}\nabla\phi\cdot\vect n_c\,\mathrm dS
  =\int_{\Omega_c}\omega\,\mathrm dV .
  \label{eq:gauss}
\end{equation}
将每个面上的面积分用面中心值近似，并定义单元平均源项
\begin{equation}
  \omega_c=\frac{1}{V_c}\int_{\Omega_c}\omega\,\mathrm dV,
\end{equation}
可得有限体积空间离散的基本形式
\begin{equation}
  \boxed{
  \sum_{f\in\mathcal F_c}
  (\nabla\phi)_f\cdot\Scf
  =\omega_cV_c .}
  \label{eq:semidiscrete-local}
\end{equation}
若 OpenFOAM 中的 \(\omega_c\) 只是单元中心值，则右端 \(\omega_cV_c\) 相当于对源项采用单元中点求积。式~\eqref{eq:semidiscrete-local}对应 Moukalled 等\cite[第 8 章，式 (8.4)]{Moukalled2016}给出的扩散项有限体积面通量和，只是这里扩散系数等于 1，且方程右端记为 \(\omega\)。

\subsection{全局面通量及守恒装配}

对每个内部面只计算一次全局定向通量
\begin{equation}
  H_f=(\nabla\phi)_f\cdot\Sf .
  \label{eq:Hdef}
\end{equation}
它在 owner 单元方程中的外向贡献是 \(+H_f\)，在 neighbour 单元方程中的外向贡献是 \(-H_f\)。因此，可将式~\eqref{eq:semidiscrete-local}写成
\begin{equation}
  \sum_{f\in\mathcal F_c^{\mathrm{int}}}\sigma_{cf}H_f
  +\sum_{b\in\mathcal F_c^{\mathrm{bnd}}}H_{cb}^{\mathrm{bnd}}
  =\omega_cV_c ,
  \label{eq:globalfluxform}
\end{equation}
其中
\(H_{cb}^{\mathrm{bnd}}=(\nabla\phi)_b\cdot\vect S_{cb}\)
按控制体外法向定义。于是未归一化残差和归一化残差分别为
\begin{align}
  Q_c(\phi)
  &:=\sum_{f\in\mathcal F_c^{\mathrm{int}}}\sigma_{cf}H_f
  +\sum_{b\in\mathcal F_c^{\mathrm{bnd}}}H_{cb}^{\mathrm{bnd}}
  -\omega_cV_c,                                      \label{eq:Qc}\\
  R_c(\phi)&:=\frac{Q_c(\phi)}{V_c}.                 \label{eq:Rc}
\end{align}
所求离散解满足
\begin{equation}
  \boxed{Q_c(\phi)=0\quad\Longleftrightarrow\quad R_c(\phi)=0,
  \qquad \forall c.}
  \label{eq:reszero}
\end{equation}
注意：\(Q_c\) 是面通量和减去体积分源项，\textbf{尚未除以} \(V_c\)；只有 \(R_c\) 才是除以体积后的残差。

内部面的守恒装配规则为
\begin{equation}
  Q_{O_f}\mathrel{+}=H_f,
  \qquad
  Q_{N_f}\mathrel{-}=H_f.
  \label{eq:assemblyflux}
\end{equation}
所以任一内部面在全域求和中严格成对抵消，这是单元中心有限体积离散的局部守恒结构。

\section{正交网格上的面通量与矩阵}

若 \(\Sf\parallel\df\)，则面法向导数可用中心差分近似：
\begin{equation}
  (\nabla\phi)_f\cdot\nf
  \approx\frac{\phi_{N_f}-\phi_{O_f}}{d_f}.
  \label{eq:orthsn}
\end{equation}
因此
\begin{equation}
  \boxed{
  H_f=D_f(\phi_{N_f}-\phi_{O_f}),
  \qquad
  D_f=\frac{|S_f|}{d_f}.}
  \label{eq:orthflux}
\end{equation}
这就是 Moukalled 等\cite[式 (8.63)]{Moukalled2016}在扩散系数为 1 时的正交网格面法向梯度公式。

原方程 \(\nabla^2\phi=\omega\)直接离散会形成负对角矩阵。为了采用常用的正定形式进行线性求解，可将连续方程两边同时乘以 \(-1\)：
\begin{equation}
  -\nabla^2\phi=-\omega .
  \label{eq:spdcontinuous}
\end{equation}
这不改变数学解。对每个内部面 \(f\)，正定矩阵 \(A\) 的装配为
\begin{equation}
\begin{array}{lll}
  A_{O_fO_f}\mathrel{+}=D_f,&
  A_{O_fN_f}\mathrel{-}=D_f,\\[0.25em]
  A_{N_fN_f}\mathrel{+}=D_f,&
  A_{N_fO_f}\mathrel{-}=D_f.
\end{array}
\label{eq:matrixassembly}
\end{equation}
内部面对应的局部块即
\begin{equation}
  D_f
  \begin{bmatrix}1&-1\\-1&1\end{bmatrix},
\end{equation}
它是对称半正定的；施加足够的 Dirichlet 边界后，全局矩阵通常成为对称正定矩阵。

\subsection{Dirichlet 边界}

设边界面 \(b\) 上规定
\(\phi_b=g_b\)，单元中心到边界面中心的距离为
\(d_{cb}=\lVert\vect x_b-\vect x_c\rVert\)。在边界中心连线与面法向一致时，原方程的外向通量为
\begin{equation}
  H_{cb}^{\mathrm{bnd}}
  \approx D_b(g_b-\phi_c),
  \qquad D_b=\frac{|S_b|}{d_{cb}}.
  \label{eq:dirflux}
\end{equation}
对正定形式~\eqref{eq:spdcontinuous}，该边界面给出
\begin{equation}
  A_{cc}\mathrel{+}=D_b,
  \qquad
  b_c\mathrel{+}=D_bg_b,
  \qquad
  b_c\text{ 的体源初值为 }-\omega_cV_c.
  \label{eq:dirassembly}
\end{equation}

OpenFOAM v14 的普通非耦合 \texttt{fixedValue} 边界通过
\(\texttt{gradientInternalCoeffs}=-\delta_b\) 和
\(\texttt{gradientBoundaryCoeffs}=\delta_bg_b\)
形成上述通量系数\cite{OFfixedValue}。需要特别区分：Moukalled 等教材在非正交边界上还推导了理论上的交叉扩散修正\cite[式 (8.81)--(8.83)]{Moukalled2016}；而 OpenFOAM v14 对普通非耦合边界将非正交修正向量置零\cite{OFsurface}。因此，若需要保持边界附近的形式精度，应尽量使边界面中心连线与边界法向接近正交，或专门设计相应边界离散。

\section{OpenFOAM v14 的非正交修正}

\subsection{教科书的几何分解}

非正交网格中，\(\Sf\) 与 \(\df\) 不再平行。教材将面积向量分解为
\begin{equation}
  \Sf=\vect E_f+\vect T_f,
  \qquad \vect E_f\parallel\df,
  \label{eq:decompose}
\end{equation}
从而
\begin{equation}
  (\nabla\phi)_f\cdot\Sf
  =|E_f|\frac{\phi_{N_f}-\phi_{O_f}}{d_f}
   +(\nabla\phi)_f\cdot\vect T_f .
  \label{eq:crossdiff}
\end{equation}
第一项可由相邻单元值隐式线性化，第二项是非正交（交叉扩散）修正。Moukalled 等\cite[式 (8.66)--(8.71)，第 8.6.5 节]{Moukalled2016}指出，交叉扩散项通常用当前梯度场显式计算，并以延迟修正方式加入代数方程右端。

单元梯度可由 Green--Gauss 公式近似：
\begin{equation}
  (\nabla\phi)_c
  \approx\frac{1}{V_c}
  \sum_{f\in\mathcal F_c}\phi_f\Scf,
  \label{eq:gggrad}
\end{equation}
其中内部面 \(\phi_f\) 可由 \(\phi_{O_f}\) 与 \(\phi_{N_f}\) 线性插值。式~\eqref{eq:gggrad}对应教材式 (8.72)--(8.76)\cite{Moukalled2016}。

\subsection{与 OpenFOAM v14 源代码一致的写法}

OpenFOAM v14 的 \texttt{correctedSnGrad} 可用统一记号准确写成
\begin{align}
  \delta_f^{\perp}
  &=\frac{1}{\max\!\left(\nf\cdot\df,\,0.05d_f\right)},
  \label{eq:deltaOF}\\
  \kf&=\nf-\delta_f^{\perp}\df,
  \label{eq:kOF}\\
  (\nabla\phi)_f\cdot\nf
  &\approx \delta_f^{\perp}
  (\phi_{N_f}-\phi_{O_f})
  +\kf\cdot(\nabla\phi)_f^{\mathrm{lin}}.
  \label{eq:snGradOF}
\end{align}
其中 \((\nabla\phi)_f^{\mathrm{lin}}\) 是单元梯度线性插值到面上的值。式~\eqref{eq:deltaOF}和~\eqref{eq:kOF}分别对应 OpenFOAM v14 的
\texttt{nonOrthDelta\allowbreak Coeffs} 与
\texttt{nonOrthCorrection\allowbreak Vectors}\cite{OFsurface}；式~\eqref{eq:snGradOF}对应
\texttt{correctedSnGrad::\allowbreak fullGradCorrection}\cite{OFcorrected}。

因此内部面全局通量为
\begin{equation}
  \boxed{
  H_f=|S_f|\left[
  \delta_f^{\perp}(\phi_{N_f}-\phi_{O_f})
  +\kf\cdot(\nabla\phi)_f^{\mathrm{lin}}
  \right].}
  \label{eq:OFflux}
\end{equation}
定义
\begin{equation}
  D_f^{\perp}=|S_f|\delta_f^{\perp},
  \qquad
  C_f=|S_f|\kf\cdot(\nabla\phi)_f^{\mathrm{lin}},
  \label{eq:DC}
\end{equation}
则
\begin{equation}
  H_f=D_f^{\perp}(\phi_{N_f}-\phi_{O_f})+C_f.
  \label{eq:HDC}
\end{equation}
在正交网格上，\(\nf\parallel\df\)，故
\(\delta_f^{\perp}=1/d_f\)、\(\kf=\vect0\)，式~\eqref{eq:OFflux}退化为式~\eqref{eq:orthflux}。公式中 \(0.05d_f\) 是 OpenFOAM v14 源代码为极端非正交网格设置的下限，不是从连续方程推导得到的数学常数\cite{OFsurface}。

对正定方程 \(-\nabla^2\phi=-\omega\)，\(D_f^{\perp}\) 形成隐式矩阵系数，\(C_f\) 作为显式非正交修正加入右端：
\begin{equation}
  b_{O_f}\mathrel{+}=C_f,
  \qquad
  b_{N_f}\mathrel{-}=C_f.
  \label{eq:Cassembly}
\end{equation}
由于 \(C_f\) 依赖当前 \(\phi\) 的梯度，非正交网格上通常需要若干次非正交修正迭代。教材第 8.10.2 节明确说明，\texttt{fvm::laplacian} 返回隐式矩阵，同时将非正交项加入源项\cite[第 8.10.2 节及代码清单 8.5--8.9]{Moukalled2016}；OpenFOAM v14 源代码也显示 \texttt{gaussLaplacianScheme} 先组装未修正矩阵，再把 \texttt{snGrad} 的显式修正加入矩阵源项\cite{OFfvmlap}。

\section{最终空间离散形式}

把上面的通量写回单元方程，原方程 \(\nabla^2\phi=\omega\) 的统一有限体积形式为
\begin{equation}
\boxed{
\begin{aligned}
  &\sum_{f\in\mathcal F_c^{\mathrm{int}}}
  \sigma_{cf}
  |S_f|\left[
  \delta_f^{\perp}(\phi_{N_f}-\phi_{O_f})
  +\kf\cdot(\nabla\phi)_f^{\mathrm{lin}}
  \right]\\
  &\hspace{4em}+
  \sum_{b\in\mathcal F_c^{\mathrm{bnd}}}
  H_{cb}^{\mathrm{bnd}}
  =\omega_cV_c .
\end{aligned}}
\label{eq:finaldiscrete}
\end{equation}
等价的残差形式为
\begin{equation}
\boxed{
  R_c(\phi)=\frac{1}{V_c}
  \left[
  \sum_{f\in\mathcal F_c^{\mathrm{int}}}\sigma_{cf}H_f
  +\sum_{b\in\mathcal F_c^{\mathrm{bnd}}}H_{cb}^{\mathrm{bnd}}
  -\omega_cV_c
  \right]=0.}
\label{eq:finalres}
\end{equation}
这就是第 4.1 问所要求的空间“半离散”方程。若网格正交，只需令
\(\delta_f^{\perp}=1/d_f\)、\(\kf=\vect0\)，即可得到中心差分形式。

\section{基于 OpenFOAM 空间离散的求解器}

\subsection{核心 C++ 程序}

以下程序面向 OpenFOAM Foundation v14。为了得到正对角的对称系统，代码求解与原方程等价的
\(-\nabla^2\phi=-\omega\)。\texttt{simpleControl} 提供稳态外迭代及非正交修正循环；空间离散方法由算例的 \texttt{system/fvSchemes} 在运行时选择。

\begin{lstlisting}[language=C++,caption={poissonFoam.C}]
#include "fvCFD.H"
#include "simpleControl.H"

int main(int argc, char *argv[])
{
    #include "setRootCase.H"
    #include "createTime.H"
    #include "createMesh.H"

    simpleControl simple(mesh);

    volScalarField phi
    (
        IOobject
        (
            "phi",
            runTime.name(),
            mesh,
            IOobject::MUST_READ,
            IOobject::AUTO_WRITE
        ),
        mesh
    );

    volScalarField omega
    (
        IOobject
        (
            "omega",
            runTime.name(),
            mesh,
            IOobject::MUST_READ,
            IOobject::NO_WRITE
        ),
        mesh
    );

    Info<< "Solving laplacian(phi) = omega\n" << endl;

    while (simple.loop(runTime))
    {
        while (simple.correctNonOrthogonal())
        {
            // Equivalent positive-definite form:
            //     -laplacian(phi) = -omega
            fvScalarMatrix phiEqn
            (
                -fvm::laplacian(phi) == -omega
            );

            phiEqn.solve();
        }

        runTime.write();
    }

    Info<< "End\n" << endl;
    return 0;
}
\end{lstlisting}

若只写与原方程同号的最简表达，也可以使用
\begin{lstlisting}[language=C++]
fvScalarMatrix phiEqn(fvm::laplacian(phi) == omega);
phiEqn.solve();
\end{lstlisting}
两者解完全相同；前一种写法更直观地对应正定 Poisson 矩阵。OpenFOAM 中
\texttt{fvc::\allowbreak laplacian(phi)} 返回已显式计算的单元场，而
\texttt{fvm::\allowbreak laplacian(phi)} 返回可组装和求解的
\texttt{fvMatrix}，因此本题应采用后者\cite[第 8.10.2 节]{Moukalled2016}。

\subsection{编译文件}

\begin{lstlisting}[language=OpenFOAM,caption={Make/files}]
poissonFoam.C

EXE = $(FOAM_USER_APPBIN)/poissonFoam
\end{lstlisting}

\begin{lstlisting}[language=OpenFOAM,caption={Make/options}]
EXE_INC = \
    -I$(LIB_SRC)/finiteVolume/lnInclude \
    -I$(LIB_SRC)/meshTools/lnInclude

EXE_LIBS = \
    -lfiniteVolume \
    -lmeshTools
\end{lstlisting}

在求解器目录中执行 \texttt{wmake} 即可编译。

\subsection{空间格式：system/fvSchemes}

\begin{lstlisting}[language=OpenFOAM,caption={本题所需的 fvSchemes 核心设置}]
gradSchemes
{
    default         Gauss linear;
    grad(phi)       Gauss linear;
}

laplacianSchemes
{
    default         none;
    laplacian(phi)  Gauss linear corrected;
}

interpolationSchemes
{
    default         linear;
}

snGradSchemes
{
    default         corrected;
}
\end{lstlisting}

这组设置逐项含义如下：
\begin{enumerate}
  \item \texttt{Gauss}：使用 Gauss 定理把单元体积分转成面通量和；OpenFOAM 用户指南将其作为 Laplacian 离散的 Gauss 形式\cite{OFGuide}。
  \item \texttt{linear}：面上所需的量采用相邻单元中心值线性插值。
  \item \texttt{corrected}：使用式~\eqref{eq:snGradOF}的显式非正交修正；其实现由 \texttt{correctedSnGrad} 源代码直接给出\cite{OFcorrected}。
\end{enumerate}
在完全正交网格上可把 \texttt{corrected} 换成 \texttt{orthogonal}；但对于一般非结构网格，直接使用 \texttt{orthogonal} 会忽略交叉扩散项。

\subsection{线性求解与非正交修正：system/fvSolution}

对乘以 \(-1\) 后的对称正定系统，可采用 PCG 与 DIC 预条件。以下参数是一组可复现实验的选择，并非原题唯一规定的数值：
\begin{lstlisting}[language=OpenFOAM,caption={fvSolution 核心设置}]
solvers
{
    phi
    {
        solver          PCG;
        preconditioner  DIC;
        tolerance       1e-10;
        relTol          0;
    }
}

SIMPLE
{
    nNonOrthogonalCorrectors  2;

    residualControl
    {
        phi             1e-10;
    }
}
\end{lstlisting}
对接近正交的网格，\texttt{nNonOrthogonalCorrectors} 可设为 0；网格非正交性明显时可使用 1--2 次修正，并通过残差及网格敏感性检查判断是否充分。OpenFOAM v14 的非正交控制类从算法字典读取该参数并执行修正循环\cite{OFnonOrthControl}。

\subsection{场文件及边界条件}

若 \(\phi\) 无量纲，则 \(\omega\) 的量纲为 \(L^{-2}\)。一个边界值场骨架为
\begin{lstlisting}[language=OpenFOAM,caption={0/phi 的边界骨架}]
dimensions      [0 0 0 0 0 0 0];
internalField   uniform 0;

boundaryField
{
    left
    {
        type        fixedValue;
        value       uniform 0;  // 用题目给定的 g 替换
    }

    right
    {
        type        fixedValue;
        value       uniform 0;
    }

    frontAndBack
    {
        type        empty;      // 二维挤出网格
    }
}
\end{lstlisting}
源项场骨架为
\begin{lstlisting}[language=OpenFOAM,caption={0/omega 的量纲与内部场}]
dimensions      [0 -2 0 0 0 0 0];
internalField   uniform 0;  // 替换为题目给定的 omega(x,y)
\end{lstlisting}
若 \(\omega\) 是坐标函数，可以在求解器中由网格单元中心 \texttt{mesh.C()} 计算，也可以先用 \texttt{setExprFields} 生成场。第 4.1 问只给出了抽象的 \(\omega\)，所以这里不擅自添加具体函数。

\section{实现核查与验证要点}

\begin{enumerate}
  \item \textbf{内部面守恒：}按式~\eqref{eq:assemblyflux}，每个内部面在 owner 和 neighbour 中的贡献符号相反；全域求和后只剩边界通量与总体源项。
  \item \textbf{总体平衡：}收敛后应满足
  \[
    \sum_{c}\omega_cV_c
    \approx
    \sum_{b\subset\partial\Omega}H_b^{\mathrm{bnd}}.
  \]
  \item \textbf{常数解检查：}当 \(\omega=0\) 且所有 Dirichlet 边界均为同一常数时，数值解应为该常数；这可同时检查符号、边界和矩阵装配。
  \item \textbf{正交退化检查：}在正交网格上，\(\kf=\vect0\)，\texttt{corrected} 与 \texttt{orthogonal} 应给出相同结果（舍入误差除外）。
  \item \textbf{网格收敛：}若给出光滑精确解和相容边界，可比较 \(L_1,L_2,L_\infty\) 误差随网格加密的变化。该内容属于后续第 4.2 问的数值验证，不应反向添加进第 4.1 问的已知条件。
\end{enumerate}

\section{结论}

第 4.1 问的核心答案是式~\eqref{eq:finaldiscrete}或等价残差式~\eqref{eq:finalres}。在 OpenFOAM v14 中，
\begin{equation*}
  \texttt{fvm::laplacian(phi)}
  \quad+\quad
  \texttt{Gauss linear corrected}
\end{equation*}
正好实现“Gauss 面通量和 + 线性面插值 + 显式非正交修正”。内部面采用 \(O_f\to N_f\) 的单一全局方向，只计算一次 \(H_f\)，然后 owner 加、neighbour 减；残差记号始终为未归一化的 \(Q_c\) 与归一化的 \(R_c=Q_c/V_c\)。

\begin{thebibliography}{99}

\bibitem{Moukalled2016}
F. Moukalled, L. Mangani, and M. Darwish,
\textit{The Finite Volume Method in Computational Fluid Dynamics: An Advanced Introduction with OpenFOAM and Matlab},
Springer, 2016, DOI: \href{https://doi.org/10.1007/978-3-319-16874-6}{10.1007/978-3-319-16874-6}.
本解答主要使用第 8 章：式 (8.4)、(8.63)--(8.77)、(8.81)--(8.83)，以及第 8.10.2 节（第 211--265 页）。

\bibitem{OFGuide}
OpenFOAM Foundation,
\textit{OpenFOAM User Guide: Numerical schemes},
\url{https://doc.cfd.direct/openfoam/user-guide/fvschemes},
访问日期：2026-08-27。

\bibitem{OFfvmlap}
OpenFOAM Foundation,
\textit{OpenFOAM v14 source: fvmLaplacian.C},
\url{https://cpp.openfoam.org/v14/fvmLaplacian_8C_source.html},
访问日期：2026-08-27。

\bibitem{OFcorrected}
OpenFOAM Foundation,
\textit{OpenFOAM v14 source: correctedSnGrad.C},
\url{https://cpp.openfoam.org/v14/correctedSnGrad_8C_source.html},
访问日期：2026-08-27。

\bibitem{OFsurface}
OpenFOAM Foundation,
\textit{OpenFOAM v14 source: surfaceInterpolation.C},
\url{https://cpp.openfoam.org/v14/surfaceInterpolation_8C_source.html},
访问日期：2026-08-27。

\bibitem{OFfixedValue}
OpenFOAM Foundation,
\textit{OpenFOAM v14 source: fixedValueFvPatchField.C},
\url{https://cpp.openfoam.org/v14/fixedValueFvPatchField_8C_source.html},
访问日期：2026-08-27。

\bibitem{OFnonOrthControl}
OpenFOAM Foundation,
\textit{OpenFOAM v14 source: nonOrthogonalSolutionControl.C},
\url{https://cpp.openfoam.org/v14/nonOrthogonalSolutionControl_8C_source.html},
访问日期：2026-08-27。

\end{thebibliography}

\end{document}
